% !TeX root = ../navier-stokes-manuscript.tex
% ============================================================
% 2.1 Vortex Stretching
% ============================================================

The energy estimate leaves one fact settled: the kinetic energy stays finite throughout the maximal smooth interval.
That control does not stop the velocity from changing more sharply across shorter distances.
Viscosity is the part of the equation we expect to stop that sharpening.
It acts alongside the three other pieces of the same motion:
\[
  \underbrace{\partial_t u}_{\text{local acceleration}}
  +
  \underbrace{(u\cdot\nabla)u}_{\text{nonlinear transport}}
  =
  \underbrace{-\nabla p}_{\text{pressure}}
  +
  \underbrace{\nu\Delta u}_{\text{viscous diffusion}}.
\]
The regularity hope is that viscous smoothing always keeps pace with the sharpening produced by nonlinear transport.
In three dimensions, vortex stretching gives the nonlinear motion a route to sharpen the flow.

\subsection{Vortex Stretching}\label{sub:ThePerfectlyStretchingVortex}

Follow a small material element whose vorticity points along one axis.
Write \(e\) for that direction and \(L(t)\) for the element's length.
If the strain is locally uniform across the element and keeps \(e\) as an extensional direction, then
\[
  S:=\frac12\bigl(\nabla u+(\nabla u)^{\mathsf T}\bigr),
  \qquad
  Se=\sigma(t)e,
  \qquad
  \sigma(t)>0,
\]
and the length grows at the fractional rate
\[
  \frac{\dot L(t)}{L(t)}
  =
  e\cdot Se
  =
  \sigma(t).
\]

The element cannot lengthen without becoming narrower.
Incompressibility preserves its material volume \(\mathcal V\), so if
\[
  A(t):=\frac{\mathcal V}{L(t)},
  \qquad
  r_{\mathrm{eff}}(t):=\sqrt{\frac{A(t)}{\pi}},
\]
then
\[
  \frac{d}{dt}(A L)=0,
  \qquad
  \frac{\dot A}{A}=-\sigma(t),
  \qquad
  \frac{\dot r_{\mathrm{eff}}}{r_{\mathrm{eff}}}
  =
  -\frac{\sigma(t)}{2}.
\]
Axial extension has forced transverse contraction in the same material element.

The curl of the momentum equation tells us what happens to the rotation along that trajectory.
With \(\omega:=\nabla\times u\),
\[
  \frac{D\omega}{Dt}
  =
  S\omega+\nu\Delta\omega,
  \qquad
  \frac{D}{Dt}:=\partial_t+u\cdot\nabla.
\]
If the vorticity is aligned with the extensional direction, \(\omega=|\omega|e\), then
\[
  S\omega=\sigma(t)\omega.
\]
Along that aligned trajectory, the full pointwise balance is
\[
  \frac12\frac{D|\omega|^2}{Dt}
  =
  \sigma(t)|\omega|^2
  +
  \nu\,\omega\cdot\Delta\omega.
\]
Whenever this right-hand side stays positive, the rotation strengthens while the effective radius contracts.

Vorticity is already part of the velocity gradient.
Indeed,
\[
  \left|
    \frac12\bigl(\nabla u-(\nabla u)^{\mathsf T}\bigr)
  \right|_{\mathrm F}
  =
  \frac{|\omega|}{\sqrt2}
  \leq
  |\nabla u|_{\mathrm F},
\]
while the global energy identity in \Cref{prop:ClassicalKineticEnergyIdentity} still gives
\[
  \norm{u(t)}_2
  \leq
  \norm{u_0}_2
  <\infty.
\]
The energy can remain finite while the rotational gradient rises.
The vortex stretches, narrows, and spins faster.
The gradients rise.

\divider{\NavierStokes{} Scaling}

At a smaller length, viscosity acts more strongly, but nonlinear transport changes with it.
The exact comparison is obtained by rescaling the whole equation.
For \(\lambda>1\), set
\[
  u_\lambda(x,t)
  :=
  \lambda u(\lambda x,\lambda^2t),
  \qquad
  p_\lambda(x,t)
  :=
  \lambda^2p(\lambda x,\lambda^2t).
\]
This produces another solution with the same viscosity at the finer length \(\lambda^{-1}\).
It compares the equation across scales; it is not a later state of the material element we just followed.

Every momentum term changes by the same factor:
\[
\begin{aligned}
  \partial_tu_\lambda(x,t)
  &=
  \lambda^3(\partial_tu)(\lambda x,\lambda^2t),
  \\
  (u_\lambda\cdot\nabla)u_\lambda(x,t)
  &=
  \lambda^3\bigl((u\cdot\nabla)u\bigr)(\lambda x,\lambda^2t),
  \\
  \nabla p_\lambda(x,t)
  &=
  \lambda^3(\nabla p)(\lambda x,\lambda^2t),
  \\
  \nu\Delta u_\lambda(x,t)
  &=
  \lambda^3\nu(\Delta u)(\lambda x,\lambda^2t).
\end{aligned}
\]
A finer scale strengthens viscous diffusion, nonlinear transport, pressure, and local acceleration together.
The equation has given viscosity no relative advantage.

\divider{Critical Height}

The balance of the equation has stayed fixed, but the quantities used to measure the velocity do not behave alike under this rescaling.
Let \(\Lambda:=(-\Delta)^{1/2}\).
The Fourier transform of the rescaled velocity is
\[
  \widehat{u_\lambda}(\xi,t)
  =
  \lambda^{-2}\widehat u(\lambda^{-1}\xi,\lambda^2t).
\]
Changing variables in the \(\dot H^s\) norm gives
\[
  \norm{u_\lambda(t)}_{\dot H^s}^2
  =
  \lambda^{2s-1}
  \norm{u(\lambda^2t)}_{\dot H^s}^2.
\]
At \(s=0\), kinetic energy loses one power of \(\lambda\).
At \(s=1\), the first-derivative energy gains one.
The scale-neutral level lies between them at \(s=\tfrac12\).
Set
\[
  H(t)
  :=
  \frac12\norm{u(t)}_{\dot H^{1/2}}^2
  =
  \frac12\norm{\Lambda^{1/2}u(t)}_2^2.
\]
Writing \(H_\lambda\) for the same quantity evaluated on \(u_\lambda\), the three measurements now sit beside one another:
\[
  \norm{u_\lambda(t)}_2^2
  =
  \lambda^{-1}\norm{u(\lambda^2t)}_2^2,
  \qquad
  H_\lambda(t)=H(\lambda^2t),
  \qquad
  \norm{\nabla u_\lambda(t)}_2^2
  =
  \lambda\norm{\nabla u(\lambda^2t)}_2^2.
\]
The finer comparison can carry less kinetic energy and a larger gradient while \(H\) keeps exactly the same size.
This scale-invariant whole-field quantity is the critical height.

\divider{Nonlinear Production versus Viscous Dissipation}

The original velocity history now tells us whether that height is rising or falling.
Keep pressure inside the nonlinear side until the global pairing is formed, and define
\[
  \mathcal P_{1/2}[u](t)
  :=
  -\operatorname{Re}
  \pair
    {\Lambda^{1/2}u(t)}
    {\Lambda^{1/2}\bigl((u(t)\cdot\nabla)u(t)+\nabla p(t)\bigr)}_{L^2}.
\]
Because \(\Lambda\) commutes with spatial derivatives and \(u\) is divergence-free, the pressure part vanishes in this global pairing.
The viscous part integrates to \(-\nu\norm{\Lambda^{3/2}u}_2^2\), leaving the exact balance
\[
  H'(t)
  +
  \nu\norm{\Lambda^{3/2}u(t)}_2^2
  =
  \mathcal P_{1/2}[u](t).
\]
The critical height rises only when the complete signed nonlinear production exceeds the simultaneous viscous loss.

Under the same rescaling,
\[
  \mathcal P_{1/2}[u_\lambda](t)
  =
  \lambda^2\mathcal P_{1/2}[u](\lambda^2t),
  \qquad
  \nu\norm{\Lambda^{3/2}u_\lambda(t)}_2^2
  =
  \lambda^2\nu\norm{\Lambda^{3/2}u(\lambda^2t)}_2^2.
\]
This is the scale match at the centre of the problem.
When the length is reduced by \(\lambda^{-1}\), nonlinear production and viscous dissipation both accelerate by \(\lambda^2\).
Passing to a finer scale does not tip their balance toward viscosity.

\Needspace{15\baselineskip}
\divider{The Heat Clock}

The same factor fixes the amount of time associated with a smaller length.
For the linear heat evolution \(v_t=\nu\Delta v\),
\[
  \widehat v(\xi,t+\delta t)
  =
  e^{-\nu|\xi|^2\delta t}\widehat v(\xi,t).
\]
Call \(r:=|\xi|^{-1}\) the length seen by a Fourier band concentrated near \(|\xi|\simeq r^{-1}\).
The multiplier changes by an order-one amount when
\[
  \nu|\xi|^2\delta t\simeq1,
\]
which gives the viscous comparison time
\[
  \tau_\nu(r)
  :=
  \frac{r^2}{\nu}.
\]
Viscosity acts throughout the interval; \(\tau_\nu(r)\) is the time required for an order-one viscous change at that spectral length.
Halving the length quarters the clock:
\[
  \tau_\nu(2^{-1}r)
  =
  \frac14\tau_\nu(r),
\]
and, more generally,
\[
  \tau_\nu(\lambda^{-1}r)
  =
  \lambda^{-2}\tau_\nu(r).
\]
This is exactly the time contraction in the full \NavierStokes{} rescaling.
The material radius in the vortex picture and the spectral length in this calculation are different measurements; an actual concentrating history has to connect them rather than assume that connection.

\divider{The Zeno Cascade}

Now repeat the scale comparison.
For the dyadic lengths
\[
  r_n:=2^{-n}r_0,
  \qquad
  \tau_n:=\frac{r_n^2}{\nu},
\]
the heat clocks form the geometric sequence
\[
  \tau_n
  =
  \frac{r_0^2}{\nu}4^{-n},
  \qquad
  \sum_{n=0}^{\infty}\tau_n
  =
  \frac{4r_0^2}{3\nu}
  <
  \infty.
\]
This finite sum is the Zeno image.
Each finer length receives less time, yet the complete sequence of comparison clocks fits below a finite ceiling.

The arithmetic does not produce a fluid cascade.
For these clocks to describe a candidate singularity, one \NavierStokes{} solution must actually pass through successively finer lengths
\[
  r_0>r_1>r_2>\cdots\longrightarrow0
\]
at times
\[
  t_0<t_1<t_2<\cdots\uparrow T_*<\infty,
  \qquad
  t_{n+1}-t_n\asymp\tau_n.
\]
The comparison constants in this relation must not deteriorate with \(n\).
The same history must also realize the connection between its material concentration and the spectral lengths measured by the heat clock.
Under those requirements, the remaining time satisfies
\[
  T_*-t_n
  \asymp
  \sum_{k=n}^{\infty}\tau_k
  =
  \frac{4r_n^2}{3\nu}.
\]
The next transition and the total time left now have the same order, \(r_n^2/\nu\), at every rung.
As \(r_n\) tends to zero, those intervals collapse with it.
The spatial cascade has become a demand on how quickly the velocity can respond in time: \(u\), then \(\partial_tu\), then \(\partial_t^2u\), and onward through the same equation.
